\documentclass{article}

\usepackage{PRIMEarxiv}
\usepackage[style=numeric-comp, backend=biber]{biblatex}
\addbibresource{references.bib}
\usepackage[T1]{fontenc}
\usepackage{amsmath}
\usepackage{booktabs}       % professional-quality tables
\usepackage{amsfonts}       % blackboard math symbols
\usepackage{nicefrac}       % compact symbols for 1/2, etc.
\usepackage{microtype}      % microtypography
\usepackage{siunitx}
\usepackage{titling}
\usepackage[none]{hyphenat}
\usepackage{ragged2e}
\sisetup{per-mode=symbol, range-units=single}
% \usepackage{parskip}
\usepackage{tabularx}
\usepackage[version=4]{mhchem}
\usepackage{fancyhdr}       % header
\usepackage{graphicx}       % graphics
\graphicspath{{images/}}    % organize your images and other figures under images/
% Unit declarations carried over from the main paper so that material pasted in
% from Balloon-Pulse-Propulsion compiles without further edits.
\DeclareSIUnit{\year}{yr}
\DeclareSIUnit{\bar}{bar} % siunitx removed \bar from the SI brochure; define it ourselves (renders identically)
\DeclareSIUnit{\dBA}{dBA} % A-weighted decibels, used for launch-noise comparisons
\DeclareSIUnit{\psf}{psf} % pounds per square foot, the conventional unit for sonic-boom overpressure
\DeclareSIUnit{\barn}{b}  % nuclear cross-section unit
\usepackage{hyperref}       % hyperlinks; loaded late so it can patch the packages above
\renewcommand{\subsubsectionautorefname}{Section}
\renewcommand{\subsectionautorefname}{Section}
\renewcommand{\sectionautorefname}{Section}

%Header
\pagestyle{fancy}

% Update your Headers here
\fancyhead[LO]{PuffSat Pulsed Propulsion}
% \fancyhead[RE]{Firstauthor and Secondauthor} % Firstauthor et al. if more than 2 - must use \documentclass[twoside]{article}

\newcommand{\subtitle}[1]{\def\thesubtitle{#1}}
% Velocity change. Always write $\deltav$ (lowercase v, matching v for velocity).
\newcommand{\deltav}{\Delta v}
% Placeholder for a citation still being added to the parent bibliography.
% Renders visibly so it cannot ship unnoticed; grep for it before release.
\newcommand{\citepending}[1]{\textbf{[cite pending: #1]}}

\pretitle{%
  \begin{center}
  \LARGE
}
\posttitle{%
  \par\vskip0.5em%
  {\large\textsl{\thesubtitle}}%
  \end{center}%
  \vskip1.5em
}

%% Title
\title{PuffSat Pulsed Propulsion}
\subtitle{A Condensed Summary With Near-Term Applications}
\author{
  Seth Katz \\
  \texttt{katzseth22202@gmail.com} \\
}

\newcolumntype{L}{>{\RaggedRight\arraybackslash}X}
\newcolumntype{C}{>{\RaggedRight\arraybackslash}c}

\begin{document}
\maketitle

\begin{abstract}
PuffSat Pulsed Propulsion is a non-nuclear kinetic analogue of Project Orion's nuclear pulse propulsion. Small, low-cost satellites (``PuffSats'') fly an Earth--Jupiter--Earth gravity-assist loop that closes in 2.2 or 3.3 years. They return on retrograde orbits at Earth-relative speeds of \qtyrange{55}{68}{\kilo\meter\per\second}. At these speeds, a PuffSat's kinetic energy per unit mass exceeds methalox's specific combustion energy by more than a hundredfold.

A suborbital rocket lifts the target vehicle above the atmosphere, and PuffSat pulses supply the remaining $\deltav$. Each arriving PuffSat disperses into a gas cloud that collides with propellant released by the target. The hot gas either strikes a pusher plate or expands through a pulsed nozzle that can thrust in any direction. We argue that each target vehicle can send more new PuffSat mass toward Jupiter than the incoming mass it consumes. Repeated across generations, this surplus would sustain exponential growth in launch capacity with no further orbital-class chemical boosters. We also sketch a variant that adds an Oberth maneuver (a velocity change at perihelion, where high speed multiplies the energy gained) and returns PuffSats at Earth-relative speeds above \SI{150}{\kilo\meter\per\second}.

Finally, we apply the architecture to orbital data centers. Data centers are now wartime targets, so we argue they belong in orbits distant enough that an attack takes days rather than minutes to arrive. AI training and coding inference tolerate the seconds of light delay to Earth that this adds. Reaching those orbits takes a large $\deltav$, and under the stated assumptions PuffSat Pulsed Propulsion can deliver mass there more cheaply than fully reusable rockets such as SpaceX Starship.
\end{abstract}

\section{Introduction}\label{sec:introduction}
\textit{PuffSat Pulsed Propulsion} is a proposed non-nuclear propulsion system that uses the kinetic energy of small, low-cost satellites to accelerate a target vehicle. Each satellite, called a PuffSat, flies an Earth--Jupiter--Earth gravity-assist loop and returns on a retrograde heliocentric orbit. The return speed relative to Earth is \qtyrange{55}{68}{\kilo\meter\per\second}. At these speeds, a PuffSat's kinetic energy per unit mass exceeds methalox's specific combustion energy by more than a hundredfold. The Jupiter flyby reverses each PuffSat's orbital direction and sends it sunward. The loop therefore closes in 2.2 or 3.3 years (two or three Earth--Jupiter synodic periods), rather than the 5.5 years of two minimum-energy legs. PuffSats depart with an Earth-relative hyperbolic excess speed of \qtyrange{12.0}{14.4}{\kilo\meter\per\second}. New Horizons had an excess speed of \SI{12.6}{\kilo\meter\per\second} and reached Jupiter in 404 days~\cite{guo2008newhorizons_design}, and Ulysses had an excess speed of \SI{11.3}{\kilo\meter\per\second} and reached it in 16 months~\cite{marsden1992ulysses_jupiter}. \autoref{sec:jupiter_gravity_assists} gives the trajectory details, and the calculations are in the companion code repository~\cite{Katz_aim_is_all_you_need_2025}.

The concept builds on Project Orion~\cite{projorion}, which proposed using nuclear explosions to deliver repeated impulses to a spacecraft, and on pellet-beam propulsion~\cite{singer1979interstellar,willis1994momentum,nordley2001interstellar,nordley2015massbeam,kare2001sailbeam,davoyan2023pelletbeam,whitmire_jackson1977laser_ramjet,matloff1979runway,kare1997buzzbomb,greason2020runway,greason2022windpellet,winglee_magbeam,landis2004particlebeam}. PuffSats reach their extreme velocities through gravity assists, which are flight-proven, rather than through nuclear explosions or through beam launchers that have not been built. A conventional suborbital rocket lifts the target vehicle above the atmosphere, and PuffSat pulses supply the remaining $\deltav$. As each PuffSat approaches, its working mass disperses into a gas cloud that collides with propellant released by the target vehicle. The collision converts relative kinetic energy into hot gas. This gas either strikes a pusher plate, transferring an impulse to the target vehicle, or enters a chamber and expands through a pulsed nozzle. The nozzle directs the exhaust to produce thrust, including thrust toward the incoming PuffSats. 

Each PuffSat must preserve its payload and electronics through the interplanetary flight despite radiation, micrometeorites, and temperature extremes. Its structural components must avoid striking the target vehicle. \autoref{sec:puffsat_satellite_architecture} sketches a satellite design addressing these requirements. The target vehicle's hardware must survive repeated pulses while capturing the arriving flow. \autoref{sec:pusher_plates} and \autoref{sec:rocket_nozzles} discuss pusher plates and pulsed nozzles.

PuffSat streams are designed to fly as tight formations, so that a target vehicle meeting the first PuffSat needs only small corrections to meet the rest. PRISMA demonstrated autonomous relative navigation to 8 cm in orbit~\cite{damico2011prisma_safe}, and the VISORS mission requires centimeter-level alignment between two CubeSats~\cite{guffanti2023autonomous}. A PuffSat formation would need similar precision across thousands of PuffSats. A formation that holds its shape can still arrive off target, so the target vehicle must measure each PuffSat's offset and correct between successive impacts. \autoref{sec:guidance_navigation} addresses these guidance, navigation, and control problems.

Launched payloads can include replacement PuffSats as well as useful cargo. If each returning generation enables a greater mass of PuffSats to be sent onto new Jupiter trajectories, the returning mass grows exponentially from cycle to cycle. Under that condition, arriving PuffSats propel both cargo and the next generation, allowing the cycle to continue without further orbital-class chemical boosters after conventional launches seed it. The initial launch expense is then spread over an increasing delivered mass. \autoref{sec:exponential_mass_growth} examines this growth condition.

Large orbital data centers provide a potential near-term market for very heavy lift~\cite{aguera2026_suncatcher,nvidia2025_starcloud,foust2026_starcloud_fcc}. Iran's 2026 strikes on terrestrial data centers~\cite{cnbc2026_aws_strike,gatech2026_aws_strike} show that data centers are military targets. China's 2007 destruction of the FY-1C weather satellite~\cite{nasa_odqn2007_fengyun,kan2007_asat} confirms that direct-ascent suborbital missiles can reach sun-synchronous orbit in minutes. In \autoref{sec:data_center_attacks}, we argue for orbits distant enough that an attack takes days to arrive, leaving time to respond. Because each cluster is co-located~\cite{aguera2026_suncatcher}, only its link to Earth slows down, and AI training and long-running coding inference can tolerate the added round-trip latency of \qtyrange{2.6}{10}{\second}~\cite{anthropic_message_batches,openai_batch_api}. Reaching them takes more $\deltav$, which favors PuffSat Pulsed Propulsion over conventional rockets. Launch must also be cheap enough for orbital computing to compete with data centers on the ground. Starcloud's chief executive puts that price near \$500 per kilogram~\cite{bennici2026_johnston}. Google's Suncatcher team sets a stricter bar, finding that at \$200 per kilogram, launch spread over a satellite's life costs about as much as a terrestrial data center's energy~\cite{aguera2026_suncatcher}. \autoref{sec:launch_cost_speculation} argues that PuffSat Pulsed Propulsion can deliver mass to these orbits for less than fully reusable rockets such as SpaceX Starship~\cite{starship}.

\autoref{sec:oberth_solar_maneuvers} sketches a variant that adds an Oberth maneuver near the Sun to the Earth--Jupiter loop, producing Earth-relative speeds above \SI{150}{\kilo\meter\per\second}. This paper complements the long-term conceptual roadmap of \textit{Aim Is All You Need}~\cite{katz2026aim} by focusing on near-term applications.

\section{PuffSat Satellite Architecture}\label{sec:puffsat_satellite_architecture}
\section{Jupiter Gravity-Assist Details}\label{sec:jupiter_gravity_assists}
\section{Pusher Plates}\label{sec:pusher_plates}
\section{Pulsed Rocket Nozzles}\label{sec:rocket_nozzles}
\section{Guidance And Navigation}\label{sec:guidance_navigation}
\section{Exponential Mass Growth}\label{sec:exponential_mass_growth}
\section{Data Centers, Anti-Satellite Weapons, And Launch Costs}\label{sec:data_center_attacks}
Current events suggest orbital data centers may be attacked by hostile states.   In March 2026, Iranian drones struck Amazon Web Services data centers in the United Arab Emirates and Bahrain, the first deliberate military attack on hyperscale cloud infrastructure~\cite{cnbc2026_aws_strike,gatech2026_aws_strike}. In January 2007, China destroyed the Fengyun-1C weather satellite at \qty{865}{\kilo\meter} with a suborbital ground-launched missile~\cite{kan2007_asat,nasa_odqn2007_fengyun}. China and Russia now both field ground-launched weapons capable of attacking satellites in low Earth orbit~\cite{dia2022_space_security}.  

Attacking critical economic choke points with cheap weapons has become a militarily prominent strategy of several global conflicts. Russian missiles and drones damaged about half of Ukraine's large substations during 2022--23~\cite{iea2024_ukraine_winter}, while Ukrainian drones have repeatedly damaged Russian refineries and ports~\cite{iea2026_omr_april}. In March 2026, attacks in the Middle East and restrictions on tanker traffic through the Strait of Hormuz reduced global oil supply by 10.1 million barrels per day~\cite{iea2026_omr_april}. 

Current orbital-data-center proposals occupy altitudes already accessible to anti-satellite weapons. Project Suncatcher~\cite{aguera2026_suncatcher} and Starcloud~\cite{foust2026_starcloud_fcc} both propose dawn--dusk sun-synchronous orbits, with Starcloud using shells between \qtyrange{600}{850}{\kilo\meter}. Fengyun-1C was intercepted at \qty{865}{\kilo\meter}. 

SpaceX has launched nearly 6,000 Starlink satellites and can build 55 per week~\cite{spacex2024_sustainability}. The director of the Space Development Agency describes the intended economics of its own proliferated constellation by noting that ``it will cost more to shoot down a single satellite than it will cost to build that single satellite''~\cite{hadley2023_tournear}. Starcloud likewise distributes satellites among many orbital shells~\cite{foust2026_starcloud_fcc}, so a single attack could not reach the entire constellation. Within each tightly spaced computing cluster, however, many expensive nodes remain accessible from a single intercept trajectory.

By contrast, orbital data centers have very expensive clusters of infrastructure in tight formations.  Their nodes must remain close enough for high-bandwidth interconnects: Suncatcher's illustrative architecture places 81 satellites within a \qty{1}{\kilo\meter} radius, separated by roughly \qtyrange{100}{200}{\meter}, for about \qty{2.3}{\mega\watt} of computing at \qty{28}{\kilo\watt} per satellite~\cite{aguera2026_suncatcher}. If each satellite presents \qty{10}{\square\meter} to a cloud of \qty{1}{\gram} pellets, approximately \qty{300}{\kilo\gram} of pellets gives one expected hit per satellite across such a disk~\cite{katz2026aim}. A reference study models a Nodong-class missile lofting \qty{500}{\kilo\gram} of pellets to \qty{600}{\kilo\meter}~\cite{wright2005_physics_space_security}. That study finds an unguided pellet weapon marginal because its cloud can miss the target by kilometers. 

Terminal homing would remove that limitation.  Hiding large data-center clusters is difficult. Starcloud's solar and cooling panels would span roughly \qty{4}{\kilo\meter} on a side~\cite{nvidia2025_starcloud}, while nearly all electrical power consumed by the data center ultimately appears as waste heat. Reflected sunlight and thermal emission therefore provide strong observables for terminal guidance, and homing systems using commercial sensors are within the capabilities of many countries~\cite{wright2005_physics_space_security}. The radiators are also physically associated with the computing hardware they cool. A seeker aimed at the strongest thermal source would consequently approach the hardware rather than a remote decoy, and disabling a satellite may require striking an area much smaller than its total projected area~\cite{wright2005_physics_space_security}. Finally, a sun-synchronous orbit passes over nearly every latitude as Earth rotates beneath it~\cite{nasa_orbit_catalog}, giving geographically dispersed launch sites repeated opportunities for an intercept.

\begin{table}[htbp]
\centering
\small
\caption{Candidate data-center sites. Uncited entries are two-body calculations that ignore plane changes. $\deltav$ is a Hohmann transfer from a \qty{300}{\kilo\meter} parking orbit, including the burn to match the site. Interceptor transit runs from a straight-up coast at \qty{14}{\kilo\meter\per\second} from \qty{200}{\kilo\meter} to a minimum-energy transfer, or to a coast at escape speed for Sun--Earth L1~\cite{katz2026aim}. The sun-synchronous transit is a ballistic climb to \qty{865}{\kilo\meter}, about \qty{8}{\minute} ignoring boost. Sun-synchronous drag assumes Suncatcher's \qty{575}{\kilo\gram} satellite~\cite{aguera2026_suncatcher} with about \qty{100}{\square\meter} of array face-on to the flow, enough for \qty{28}{\kilo\watt} at 20~percent cell efficiency. Earth--Moon L4/L5 eclipses assume the point crosses Earth's shadow at the Moon's synodic rate.}
\label{tab:data_center_sites}
\begin{tabularx}{\textwidth}{@{}>{\raggedright\arraybackslash}p{0.17\textwidth}*{5}{>{\raggedright\arraybackslash}X}@{}}
\toprule
Site & Round-trip latency & Interceptor transit & $\deltav$ from \qty{300}{\kilo\meter} & Station keeping & Eclipses \\
\midrule
Sun-synchronous, \qtyrange{600}{850}{\kilo\meter} & under \qty{0.01}{\second} & Minutes & \qtyrange{0.17}{0.30}{\kilo\meter\per\second} & \qtyrange{3}{50}{\meter\per\second} a year against drag~\cite{vallado_atmosexp}, plus formation keeping~\cite{aguera2026_suncatcher} & Seasonal, up to \qty{10}{\minute} per \qty{98}{\minute} orbit~\cite{nasa_hinode} \\
\addlinespace
High Earth orbit, \qtyrange{20000}{80000}{\kilo\meter} & \qtyrange{0.13}{0.53}{\second} & \qty{0.5}{\hour} to \qty{14}{\hour} & \qtyrange{3.45}{4.14}{\kilo\meter\per\second} & Little outside a GEO slot, which takes \qty{45}{\meter\per\second} a year~\cite{oleson1995_geo_nssk} & Two seasons a year, up to \qty{72}{\minute} a day at GEO~\cite{noaa_goes_eclipse} \\
\addlinespace
Earth--Moon L4/L5 & \qty{2.6}{\second} & Half a day to 5 days & \qty{3.94}{\kilo\meter\per\second}, \qty{0.83}{\kilo\meter\per\second} of it on arrival & Little, since the points are stable~\cite{nasa_lagrange_points} & Two seasons a year, up to \qty{2.7}{\hour} in Earth's umbra \\
\addlinespace
Sun--Earth L1/L2 & \qty{10}{\second} & 2 to 16 days & \qty{3.17}{\kilo\meter\per\second} to L1~\cite{song2021_sel_launch} & About \qty{1}{\meter\per\second} a year at L1~\cite{roberts2011_l1_ace_soho_wind} & None at L1~\cite{nasa_lagrange_points} \\
\bottomrule
\end{tabularx}
\end{table}

Moving farther from Earth does not make an orbit unpredictable. a published ephemeris remains predictable at any altitude. It does, however, increase both interceptor flight time and the launch energy required. \autoref{tab:data_center_sites} compares four candidate locations. A direct-ascent interceptor can reach sun-synchronous orbit within minutes, leaving little time to detect the launch and respond. The Fengyun-1C interceptor was launched from a road-mobile vehicle~\cite{kan2007_asat}. Reaching a Lagrange point requires nearly escape velocity: about \qty{11.0}{\kilo\meter\per\second} from \qty{200}{\kilo\meter}, compared with about \qty{3.4}{\kilo\meter\per\second} for a ballistic climb to \qty{650}{\kilo\meter}. With \qty{250}{\second} solid-propellant specific impulse, the corresponding ideal mass ratio is 22 times larger~\cite{katz2026aim}.

Higher interceptor energy can shorten the trip but cannot eliminate the warning time created by distance. At \qty{14}{\kilo\meter\per\second}, an interceptor reaches Earth--Moon L4/L5 in roughly half a day and Sun--Earth L1 in roughly two days. That interval permits responses unavailable against a direct-ascent LEO attack.  A cluster can disperse, and defensive interceptors can be launched before the attacker arrives. Sun--Earth L1 therefore changes the problem from a minutes-long direct-ascent engagement to one with days of warning.

The communications penalty is modest for some computing workloads. Round-trip light time from Earth is \qty{2.6}{\second} at the Earth--Moon Lagrange points and about \qty{10}{\second} at Sun--Earth L1/L2. Batch inference services already permit turnaround times of up to \qty{24}{\hour}~\cite{anthropic_message_batches,openai_batch_api}. Distributed training is also relatively insensitive to Earth latency because the high-rate communication occurs among machines inside the cluster. Earth supplies training data and receives checkpoints rather than participating in each training step.
 
Carefully arranged Lagrange point orbits can have near continous sunlight, whereas high Earth orbits undergo eclipse seasons that require batteries or redundant satellites that take over when they are in Earth's shadow. L1 is dynamically unstable, but spacecraft there have required only about \qty{1}{\meter\per\second} of station-keeping $\deltav$ per year. Transfer to L1 also requires \qtyrange{0.3}{1.0}{\kilo\meter\per\second} less $\deltav$ than transfer to a circular orbit at \qtyrange{20000}{80000}{\kilo\meter}.   The Lagrange points are also natural places to put military defeneses that can defend all data centers at that location.   High orbits likely need to deliberately cluster around a military defense.

Radiation is the principal disadvantage. The lower part of the \qtyrange{20000}{80000}{\kilo\meter} band lies within the outer Van Allen belt, while the Lagrange points lie beyond both the belts and most of Earth's magnetospheric shielding~\cite{li2019_vanallen}. The steady galactic-cosmic-ray dose there is about 12~rad per year~\cite{zeitlin2013_msl_rad}, compared with the 150~rad(Si) per year that Suncatcher expects in its orbit behind \qty{10}{\milli\meter} of aluminum~\cite{aguera2026_suncatcher}. High orbits lose much of the magnetosphere's protection against solar particle events, requiring shielding and fault-tolerant computing to survive acute radiation events. Suncatcher's proton testing ran a commercial tensor processing unit to 15~krad(Si) without hard failures, although its high-bandwidth memory showed irregularities after 2~krad(Si)~\cite{aguera2026_suncatcher}. These results suggest that radiation imposes a substantial engineering requirement, but one that can be addressed directly through shielding, component selection, and redundancy.

Both PuffSats and Starship~\cite{starship} could in principle deliver payloads to sun-synchronous and high orbits. SpaceX's IPO underwriters have published cost projections for Starship, although underwriting forecasts have incentives toward optimism~\cite{michaely1999_underwriter}. Goldman Sachs projects \$183 per \si{\kilo\gram} to low Earth orbit by 2030~\cite{pethokoukis2026_goldman}. Morgan Stanley projects about \$500 per \si{\kilo\gram} by 2030 and less than \$150 per \si{\kilo\gram} by 2040, assuming reuse of both stages~\cite{investing2026_ms_spacex}. Full second-stage reuse requires a reusable thermal-protection system. A commercial space executive and two former NASA thermal-protection engineers, including one of the originators of PICA~\cite{nasa_spinoff2013_pica}, have argued that Starship's current tile architecture will keep costs well above \$100--250 per \si{\kilo\gram}~\cite{miller2026_tps}. SpaceX's registration statement identifies a durable reusable heat shield as a material technical challenge~\cite{spacex2026_s1a}. Its chief executive described the heat shield as Starship's largest remaining challenge in 2024~\cite{musk2024_heatshield} and said in August 2026 that he would consider the problem solved, before that flight's ship had undergone close inspection~\cite{spacex2026_q2_call}.

Deliveries beyond low Earth orbit require additional tanker launches. Sun--Earth L1 requires roughly \qtyrange{3.2}{4.0}{\kilo\meter\per\second} of additional $\deltav$, supplied by propellant transferred in orbit. At \qty{380}{\second} specific impulse, an \qty{85}{\tonne} ship~\cite{gunter_starship} delivering \qty{100}{\tonne} consumes \qtyrange{252}{356}{\tonne} of propellant, corresponding to roughly two to four tanker flights. Each kilogram delivered to L1 therefore requires about 3.5 to 4.6 kilograms launched to low Earth orbit~\cite{katz2026aim}. Applying that multiplier to the published forecasts raises Goldman's \$183 per \si{\kilo\gram} to roughly \$640--840 per \si{\kilo\gram} at L1 and Morgan Stanley's 2030 estimate to roughly \$1750--2300 per \si{\kilo\gram}.

These estimates are lower bounds. They omit the operational cost of orbital propellant transfer, which SpaceX also identifies as a material challenge~\cite{spacex2026_s1a}, and they do not account for returning the ship from L1. An L1 return reaches the atmosphere at about \qty{11.1}{\kilo\meter\per\second}, compared with roughly \qty{8.0}{\kilo\meter\per\second} from low Earth orbit. With peak heating scaling approximately as the cube of entry velocity, the corresponding peak heat flux is about 2.7 times greater~\cite{katz2026aim}. Starcloud's chief executive estimates that orbital computing becomes competitive near \$500 per \si{\kilo\gram}~\cite{bennici2026_johnston}, while Google's Suncatcher analysis uses a target of \$200 per \si{\kilo\gram}~\cite{aguera2026_suncatcher}. The published Starship cost projections exceed both thresholds when extended to the more distant sites considered here. \autoref{sec:launch_cost_speculation} argues that PuffSat Pulsed Propulsion keeps achieves the competitive launch cost target given by Starcloud and may reach it for the more stringest Suncatcher cost target.

\section{Launch Cost Speculation}\label{sec:launch_cost_speculation}
\section{Leveraging Oberth Maneuvers Near The Sun}\label{sec:oberth_solar_maneuvers}
\section{Conclusion}\label{sec:conclusion}


\appendix
\bigskip
\begingroup
  \centering
  \Large\bfseries Appendix\par
\endgroup

\section{AI Use Disclosure}\label{sec:transparency}
The underlying ideas, physical arguments, and proposed architecture in this paper and in \textit{Aim Is All You Need}~\cite{katz2026aim} are my own. The architecture began as hard-science-fiction worldbuilding, which I developed in dialogue with AI tools. I used Claude Code~\cite{claude_code}, Codex~\cite{codex}, Grok~\cite{grok}, Gemini~\cite{gemini}, and Microsoft Copilot~\cite{copilot} for formatting, coding, math and physics derivations and sanity checks, image creation, bibliography, and proofreading support. I know that's not enough, which is why I'm seeking advice and a sanity check on whether this is worth pursuing.

\section{Copyright}\label{sec:copyright}
\small
\textcopyright\ Seth Katz. All rights reserved.


%Bibliography
\printbibliography


\end{document}
